\section{Conclusion}
\label{sec:conclusion}

This report tested whether local spatial context helps a lightweight VAE
select informative MSI peaks, and which part of a learned pipeline determines
peak quality. Under matched budgets and predeclared rules, on three public
collections, we reached the following conclusions.
\begin{itemize}
  \item Ranking peaks by correlation with a label-free GMM segmentation
        outperformed gradient attribution through the VAE on every dataset.
        This was statistically supported and stable across training seeds.
  \item Neighbourhood input to the encoder and smoothing of the segmentation
        gave no consistent, spatially specific gain. Spatially aware ranking
        statistics were strongest on GBM and RCC and within 0.01 of the best
        on CAC.
  \item A linear PCA segmentation performed comparably to the VAE, and
        reconstruction error did not track peak quality.
  \item msiPL reproduced. S\textsuperscript{3}PL's published GBM score did not.
        Weight initialisation dominated seed sensitivity in two of three
        selected sections, while data order also mattered in one.
\end{itemize}

The contribution is a controlled evaluation, not a new peak picker. It shows
how spatially aware rankings performed in these datasets, and how
strongly conclusions depend on fair budgets, model-free references and seed
replication.

Future work should first:
\begin{itemize}
  \item validate a label-free choice of the number of segments, which BIC
        could not provide;
  \item test a predeclared combination of segmentation correlation and
        Moran's~$I$ on new data;
  \item evaluate spatial models that operate on graphs or whole ion images
        rather than fixed windows;
  \item test whether a model fitted on one section transfers to another.
\end{itemize}
